% \subsection{Relationship to other notions}
% To understand how merge-width compares to previously studied notions, 
% we show the following statements, each of which is easy to prove.
% The main message (made evident by the proofs) is that merge-sequences of bounded width can be 
% seen as a common generalization of the decompositions underlying treewidth, clique-width, 
% twin-width, and bounded expansion classes.

% Note that analogous statements are known to hold 
% for flip-width instead of merge-width, so the statements alone do not explicitly demonstrate 
% the advantage of merge-width over flip-width. However, the statements still
% demonstrate that merge-width (like flip-width) is a sensible extension 
% of the notions studied in Sparsity and in Twin-width.
% The  advantage of merge-width over flip-width is made apparent by our main 
% result, \Cref{thm:main}, discussed in \Cref{sec:main-res}.


% \begin{theorem}For every graph $G$,
%   $\mw_\infty(G)\le \cw(G)\le O(\mw_\infty(G))$.
% \end{theorem}
% It follows that a graph class $\CC$ has bounded treewidth if and only if 
%  $\mw_\infty(\CC)<\infty$, and $\CC$ is \emph{weakly sparse},
% that is,  excludes some biclique $K_{t,t}$ as a subgraph.

% \begin{theorem}For every graph $G$,
%   $\mw_1(G)\le \deg(G)+2$.
%   Conversely, a graph class $\CC$ has bounded degeneracy if and only if 
%    $\mw_1(\CC)<\infty$, and $\CC$ is weakly sparse.
% \end{theorem}


% \begin{theorem}
% Classes of bounded merge-width include all classes of bounded expansion, and all classes of bounded twin-width.
% Moreover:
% \begin{itemize}
%   \item a graph class  has bounded expansion if and only if it has bounded merge-width and is weakly sparse,
%   \item a class of graphs has bounded twin-width if and only if every graph in the class 
%   can be equipped with a total order, obtaining a class of ordered graphs 
%   of bounded merge-width.
% \end{itemize}
% \end{theorem}
% \noindent For the last item, note that the notion of merge-width extends easily to structures equipped with 
% one or multiple binary relations, such as ordered graphs -- graphs equipped with a total order on the vertex set.

% The next result shows that bounded merge-width implies bounded flip-width.
% We conjecture that the converse holds as well.

% \begin{theorem}
%   Every class of bounded merge-width has bounded flip-width.
%   \end{theorem}

% The following result states that classes of bounded merge-width are preserved by
% so-called first-order, one-dimensional interpretations.
% For a graph  $G$ and first-order formula  $\phi(x,y)$, 
%   we construct  
%   a graph $\phi(G)$
%   with vertices $V(G)$ and edges $\setof{uv}{u,v\in V(G), G\models\phi(u,v)}$.
% \begin{theorem}
%   Let $\phi(x,y)$ be a first-order formula and 
%    $\CC$ be a graph class of bounded merge-width.
%   Then the class $\setof{\phi(G)}{G\in\CC}$ has bounded merge-width.
% \end{theorem}


\section{Case studies}\label{sec:cases}
We now review the relationship between
merge-width and previously studied graph parameters.
First, it is not difficult to check that 
$\mw_\infty(G)$ is functionally equivalent to the clique-width of $G$. 
\begin{theorem}\label{thm:cw}
  A graph class $\CC$ has bounded clique-width if and only if $\mw_\infty(\CC)<\infty$.
%   More precisely, for every graph $G$,
%   $$\cw(G)\le \mw_\infty(G)\le 2\cw(G).$$
\end{theorem}


    It is not difficult and quite instructive to prove the result directly,
    by translating between construction sequences and clique-width expressions, and vice-versa.
    Such a proof would provide explicit (linear) bounds in each direction.
    To be more succinct, we derive the statement without explicit bounds, by deriving it from results proved later in \Cref{sec:cases}.

\begin{proof}[Proof sketch]

    For the left-to-right implication,
    let $G$ be a graph of clique-width $k$.
    By \cite[Thm. 1]{tww6}, $G$ has a contraction sequence $\cal P_1,\ldots,\cal P_n$, such that at every 
    step $i\in\set{1,\ldots,n}$ of the sequence,
    every connected component of the red graph of $\cal P_i$ contains 
    at most $k'$ parts, for some constant $k'\in\N$ depending on $k$.
    By transforming the contraction sequence into a merge sequence as in the proof of \Cref{lem:tww} below, by \eqref{eq:tww-mw} we derive that the resulting construction sequence has radius-$\infty$ width at most $k'$. Thus, $\mw_\infty(G)\le k'$.
    This proves the left-to-right implication.

    Conversely, suppose that $\mw_\infty(\CC)<\infty$.
    It follows from \Cref{thm:fw-cases} below that $\fw_\infty(\CC)<\infty$.
    Therefore, $\CC$ has bounded clique-width, by \cite[Thm. II.6]{flip-width}.
\end{proof}

% \begin{proof}[Proof sketch]\newcommand{\nlc}{\text{NLC}}
%     We use the notion of NLC-width \cite{wanke}, denoted $\nlc(G)$, which satisfies \cite{gurski-wanke}
%     $\nlc(G)\le \cw(G)\le 2\nlc(G)$.
%     Fix a construction sequence for $G$ witnessing \(\mw_\infty(G) \le k\).
% At each stage of the construction sequence, 
% with current partition $\cal P$ and set of resolved pairs $R\subset {V\choose 2}$,
% we maintain a \(k\)-expression whose \(\le k\) colors correspond to the \(\le k\) parts of the connected component.

%    In a merge-width sequence witnessing 
%    every connected component in the resolved graph consists of at most \(k\) parts.


%    Starting with the partition into singletons, we traverse the merge sequence and maintain for every connected component of the resolved graph
%    a \(k\)-expression whose \(\le k\) colors correspond to the \(\le k\) parts of the connected component.
%    Whenever edges between different connected components are resolved,
%    we first take the disjoint union of the corresponding \(k\)-expressions,
%    and then add the edges between the corresponding color classes.
% \end{proof}

%NLC width operations:
%create vertex with color i
%take disjoint union of two colored graphs G,H and arbitrarily add bicliques between any color class in G and color class in H
% recolor

\noindent Thus, in a sense, the finite-radius merge-width parameters are local variants of clique-width.




% Let $G$ be a graph, $\cal P$ be a partition of $V(G)$. For $a\in V$, let $\cal P(a)$ denote the part of $\cal P$ containing $a$.
% A \emph{default graph} on $\cal P$
% is a graph with vertex set $\cal P$, which may have self-loops, that is, $E(\Delta)\subset\setof{\set{A,B}}{A,B\in\cal P}$.
% The \emph{error graph} of $(G,\Delta)$
% is the graph with vertices $V(G)$ 
% and edge set $$\textit{error}(G,\Delta):=\setof{\set{a,b}\in {V(G)\choose 2}}{\set{a,b}\in E(G)\iff
% \set{\cal P(a),\cal P(b)}\notin E(\Delta)}.$$



% Notice that 
% if $\Pi$ is a fast merge sequence for $G$ as above, then replacing $R_i$ with  
% $$R_i'=\bigcup_{j\in [i]}\textit{error}(G,\Delta_j),$$
% for each $i\in[m]$,
% also yields a fast merge sequence $\Pi'$,
% whose radius-$r$ width is not larger than that of $\Pi$, for all $r\in\N\cup\set{\infty}$.
% It is 
% sometimes convenient to allow for larger sets $R_i$.






\subsection{Merge-width and twin-width}\label{sec:tww}
We first discuss the relationship of merge-width and twin-width,
whose definition we recall now.
Recall that a \emph{contraction sequence} of a graph $G$ is a sequence of merge operations, which starts with the partition of $V(G)$ into singletons, and ends with the partition with one part.
Two sets $A,B\subset V(G)$ of vertices are \emph{homogeneous} if $AB\subset E(G)$ or $AB\cap E(G)=\emptyset$.
The \emph{red graph} of a partition $\cal P$
is the graph with vertices $\cal P$ and edges connecting pairs $\set{A,B}\in{\cal P\choose 2}$ 
which are not homogeneous. The twin-width of a graph $G$ is the minimum number $d$ 
such that $G$ admits a contraction sequence such that at every step, the red graph of the current partition has maximum degree at most $d$.


\twwintro*
% \begin{theorem}\label{thm:tww}
%   Every graph class of bounded twin-width has bounded merge-width.
% \end{theorem}
\Cref{thm:tww} follows immediately from the next lemma.
\begin{lemma}\label{lem:tww}
  Fox every $r,d\in \N$ and graph $G$ of twin-width $d$, we have 
  $$\mw_r(G)\le 2+d+\ldots+d^r.$$
\end{lemma}
\begin{proof}
Fix a contraction sequence $\cal P_1,\ldots,\cal P_n$ of $G$, such that
the red graph of each partition $\cal P_t$ has maximum degree at most $d$.
For a vertex \(v\), denote by \(\cal P_t(v)\) the part of \(\cal P_t\) containing \(v\).
For $t\in[n]$, let $R_t$ consist of those pairs $ab\in {V\choose 2}$ such that $\cal P_t(a),\cal P_t(b)$ are either equal, or not homogeneous in $G$.
Then 
  $(\cal P_1,R_1),\ldots,(\cal P_n,R_n),$
  is a merge sequence of $G$. We bound its radius-$r$ width, for fixed $r\in\N$.

  For every $t\in[n]$, part $A\in\cal P_t$ and vertex $a\in A$, 
  we have that 
  \begin{align}\label{eq:tww-mw}
  B^r_{R_t}(a)\subset \bigcup B^r_{\cal P_t}(A),  
  \end{align}
  where $B^r_{R_t}(a)$ is the ball of radius $r$ around $a$ in the graph $(V,R_t)$, 
  and $B^r_{\cal P_t}(A)$ is the ball of radius $r$ around $A$ in the red graph of $\cal P_t$.
  In particular, $B^r_{R_t}(a)$ intersects at most $|B^r_{\cal P_t}(A)|$ parts of $\cal P_t$.
  As degree in the red graph is at most \(d\), the radius-$r$ width of $(\cal P_t,R_t)$ is at most $|B^r_{\cal P_t}(A)|\le 1+d+\ldots+d^r$.
  Since $\cal P_{t-1}$ is a refinement of $\cal P_t$ with exactly one more part,
  it follows that 
  the radius-$r$ width of $(\cal P_{t-1},R_t)$ is at most $2+d+\ldots+d^r$.
Consequently,
$\mw_r(G)\le 2+d+\ldots+d^r.$
\end{proof}

% \begin{proof}[Proof of Theorem~\ref{thm:tww}]
% A contraction sequence induces a merge sequence,
% in which just before merging any two parts $A,B\in\cal P$,
% we first resolve both $A$ and $B$ 
% with each part $C\in \cal P$ 
% such that $A\cup B$ is not homogeneous towards $C$.
% In the resulting merge sequence, at every time immediately following a merge operation, 
% a pair $\set{a,b}\in AB$, for $A,B\in\cal P$, is resolved 
% if and only if $A$ and $B$ are non-homogeneous in $G$.
% % It follows that for any two parts $A,B$ and vertices $a\in A,b\in B$,
% % the distance between $A$ and $B$ in the red graph of $\cal P$ 
% % is upper bounded by the distance between $a$ and $b$ in $(V,R)$.
% Therefore, for every $r\in\N$,
%  and every part $A\in\cal P$ and vertex $a\in V(G)$, 
%  we have that $$B^r_R(a)\subset \bigcup B^r_\cal P(A),$$
%  where $B^r_R(a)$ is the ball of radius $r$ around $a$ in the graph $(V,R)$ formed by currently resolved pairs, 
%  and $B^r_\cal P(A)$ is the ball of radius $r$ around $A$ in the red graph of $\cal P$.
%  In particular, $B^r_R(a)$ intersects at most $|B^r_\cal P(A)|$ parts of $\cal P$.
%  In the resolve steps which do not immediately follow the merge of two parts,
%  this upper bound is replaced by $|B^r_\cal P(A)|+1$.
%  As $|B^r_\cal P(A)|\le 1+d+\ldots+d^r$,
%  where $d$ is the maximum degree of the red graph of $\cal P$, 
% it follows that for all $r\in\N$, $$\mw_r(G)\le 2+\tww(G)+\cdots+\tww(G)^r\le O(\tww(G)^r).$$
% \Cref{thm:tww} follows.
% \end{proof}

\subsection{Merge-width and Sparsity}\label{sec:sparsity}
The converse to \Cref{thm:tww} is false. 
For instance,  the class 
of graphs of maximum degree bounded by $d$ (a class which has unbounded twin-width already for $d=3$) 
has bounded merge-width.
Indeed, there is a trivial merge sequence for such a graph $G$: 
 $$\bigl(\textit{partition into singletons},\emptyset\bigr),\quad \bigl(\textit{partition with one part},E(G)\bigr).$$
The radius-$r$ width of this sequence is 
the maximum size of a radius-$r$ ball in $G$,
which is at most $1+d+\cdots+d^r$, 
so 
$\mw_r(G)\le 1+d+\cdots+d^r$.


% Indeed, a merge sequence for such a graph $G$ may proceed by first resolving all 
%  pairs $\set{u},\set{v}$ of singleton parts with $uv\in E(G)$, thus 
%  revealing all the edges of $G$,
%  only after which all singleton parts are merged -- in an arbitrary sequence of merge steps -- into one single part $V(G)$;  finally, that part 
% is resolved with itself, by revealing all the non-edges of $G$.
% Apart from the last step, at any time in the sequence, every vertex $v$ reaches at most $1+d+\cdots+d^r$ vertices -- and therefore, also parts -- by a path of length at most $r$ consisting of either edges or non-edges; in particular, $\mw_r(G)\le 1+d+\cdots+d^r$.

If the above merge sequence for graphs of bounded degree seems uninsightful,
this reflects the fact that graphs of bounded degree are trivial from the perspective of Sparsity theory.
The next result is slightly more 
illuminating.
The proof is already given in \Cref{ex:degeneracy}.


%Recall that a graph $G$ is \emph{$d$-degenerate}, for $d\in\N$, if and only if there is a total order ${\le}$ on the vertices of $G$ such that every vertex $v$ has at most $d$ neighbors before it.


\begin{theorem}\label{thm:deg}
  If $G$ is a $d$-degenerate graph then 
  $\mw_1(G)\le d+2.$
\end{theorem}

%\begin{proof}
%  Fix a $d$-degenerate graph $G$ and a total order ${\le}$ witnessings its $d$-degeneracy.
%For $t\in[n]$, let $\cal P_t$ be the partition
%  in which the $t$ largest elements with respect to $\le$ form a single part $L_t$,
%  and each of the remaining vertices $v\in V(G)-L_t$ forms a singleton part $\set{v}$.
%  Let $R_t\subset E(G)$ consist of those edges of $G$ which have at least one endpoint in $L_t$.
%
%  It is straightforward to verify that
%  $(\cal P_1,R_1),\ldots,(\cal P_n,R_n)$
%  is a merge sequence for $G$.
%  Moreover, for $t\in[n]$, the radius-$1$  width of $(\cal P_t,R_t)$, is at most $d+1$.
%As $\cal P_{t-1}$ has exactly one part more than $\cal P_t$, it follows that the radius-$1$ width of $(\cal P_{t-1},R_t)$, is at most $d+2$.
%%
%%
%% Fix a $d$-degenerate graph $G$ and a total order ${\le}$ witnessings its $d$-degeneracy.
%% A merge sequence for the graph $G$
%% proceeds in rounds, where
%% at the beginning $i$th round,
%% the current partition has one part $P$
%% consisting of the $i-1$ greatest
%% vertices in the order ${\le}$
%% (in the first round $P$ is empty, which we allow for ease of description),
%% while the remaining vertices are partitioned into singletons;
%% the resolved pairs are exactly the edges of $G$ with at least one endpoint in $P$.
%% Specifically, in round $i$,
%% we first resolve the part $\set{v}$,
%% where $v$ is the $i$th greatest vertex with respect to ${\le}$,
%% with each of the parts
%% $\set{w}$  such that $w<v$ and $vw\in E(G)$, revealing all  edges incident to $v$ that haven't been resolved yet;
%% afterwards, we merge $\set{v}$ into the  part $P$. Once we are left with one part only, we resolve it with itself, revealing all the non-edges of $G$.
%% This merge sequence witnesses that $\mw_1(G)\le d$, proving \Cref{thm:deg}.
%\end{proof}

Next, we prove:
\beintro*

The simple construction from \Cref{thm:deg} does not trivially generalize to bound the radius-$r$ merge-width 
of graphs in classes of bounded expansion,
as it appears that to bound $\mw_r(G)$ for $r>1$, a more involved construction 
is needed, which we now present.

\medskip

We first recall the definition of classes of bounded expansion for completeness only, as we will not use the original definition, only the characterization given in the upcoming \Cref{fact:be}.
By definition,  a graph class $\CC$ has \emph{bounded expansion} if and only if for every $r\in\N$,
there is some $k\in\R$ such that for every graph $G$
which can be obtained from some graph in $\CC$ by first removing vertices and edges, and then contracting 
some vertex subsets of radius ${\le r}$ to single vertices, 
we have that $|E(G)|\le k\cdot |V(G)|$.
Classes of bounded expansion include every class of bounded maximum degree,
the class of planar graphs, and every graph class that excludes some graph  as a  minor, or as a topological minor.

%\[
%   -----------------------
%\]

Fix a graph $G$, number $r\in\N$, and total order $\le$ on $V(G)$. 
A vertex \(u\) is \emph{weakly \(r\)-reachable} from a vertex \(v\) if there is a path of length at most \(r\) from \(v\) 
to \(u\) and all vertices on that path are at least as large as \(u\).
Let \(\wreach_r(G,\le,v)\) be the set of vertices that are weakly \(r\)-reachable from \(v\) with respect to the ordering.
Then the weak \(r\)-coloring number of \(G\) is defined as
\[
\wcol_r(G) \coloneqq \min_{\le} \max_{v \in V(G)} |\wreach_r(G,\le,v)|,
\]
where the minimum ranges over all total orders $\le$ of $V$.

\begin{lemma}[\cite{zhu2009colouring}]\label{fact:be}
A graph class $\CC$ has bounded expansion if and only if $\wcol_r(\CC)<\infty$, for all~${r\in\N}$.
\end{lemma}


%\[
%   -----------------------
%\]


%We now define a variant of generalized coloring numbers.
%Fix a graph $G$, number $r\in\N\cup\set{\infty}$, and total order $\le$ on $V(G)$.
%Define the \emph{medium $r$-reachability number}\sz{other name} of $(G,\le)$ as the maximum,
%over all $u<v$ in $V(G)$, of the number of vertices $w\le u$ reachable
%from $v$ by a path of length ${\le r}$ whose inner vertices are strictly greater than $u$.
%Let $\mreach_r(G)$ denote the least medium $r$-reachability number of $(G,\le)$,
%over all total orders $\le$ on $V(G)$.
%
%
%Let $\scol_r$ and $\wcol_r$
%are, respectively, denote the \emph{strong} and \emph{weak} $r$-coloring numbers.
%The following relations follow immediately from the definitions,
%which we omit here (see the introduction for a definition of the weak $r$-coloring numbers).

%\begin{lemma}\label{lem:cols}
%For every graph   $G$,
%$$\scol_r(G)\le \mreach_r(G)\le \wcol_r(G).$$
%\end{lemma}


%The following fact is an immediate consequence of \Cref{lem:cols} and known results relating bounded expansion
%with generalized coloring numbers~\cite{zhu2009colouring}.

% For vertices $v,w\in V(G)$ with $w\le v$, we say that $w$ is \emph{weakly-$r$ reachable from $v$} if $w$ and $v$ are connected by a path $\pi$ of length at most $r$, such that $u\ge w$ for every vertex $u$ of $\pi$.
% The \emph{$r$-weak coloring number} of a graph $G$, denoted $\wcol_r(G)$, is the least number $k$ 
% such that there is a total order ${\le}$ 
% of the vertices of $G$ such that 
% $|\wreach^r_{G,\le}(v)|\le k$ holds
% for all $v\in V(G)$.





%See \Cref{rem:isolations} for a perspective viewing the total orders in the definition of weak coloring numbers as a special form of merge sequences.
\Cref{thm:be} follows immediately from the next lemma,
which 
shows that the 
 radius-$r$ merge-width 
is bounded in terms of the weak $(r+1)$-coloring number. 


\begin{lemma}\label{lem:wcol}
  For every $k,r\in\N$ and graph $G$ with $\wcol_{r+1}(G)\le k$, we have
  $$\mw_r(G)\le 3\cdot 2^{k} .$$
\end{lemma}
In particular, we show that $\mw_2(G)$  is bounded in terms of $\wcol_3(G)$,
and we do not know whether it is bounded in terms of $\wcol_2(G)$.

\begin{proof}Fix $r\in\N$, and a graph $G$ with vertex set $V$.
  % For two vertices $a,b$ of a graph $G$,
  % and a set $S\subset V(G)$, say that $a$ and $b$
  % have the same \emph{atomic types over $S$} if 
  % either $a=b\in S$, or $a,b\notin S$ and $N(a)\cap S=N(b)\cap S$.
  % The \emph{atomic type} of a vertex $a$ \emph{over} $S$ 
  % is the equivalence class of $a$ with respect to the above equivalence relation.
  % and is therefore uniquely determined by the 
  % pair $(\setof{u\in U}{u=a},\setof{u\in U}{E(u,a)})$ of subsets of $U$.
Let $\le$ be an ordering of the vertex set $V$ of $G$ 
such that for every vertex \(v\), the weak $r$-reachability set $\wreach_{r+1}(G,\le,v)$ has size at most $k$.

For $t=1,\ldots,n$, define the following.
\begin{itemize}
  \item Let $L_t$ comprise the $t$ largest elements of $V$ with respect to $\le$, and let $S_t\coloneqq V-L_t$.
  \item Let  $\cal P_t$ be the partition of $V$ into \emph{atomic types}
  over the set $S_t$. That is, $\cal P_t$ partitions $S_t$ into singletons, and partitions vertices  $v\in L_t$ according to their neighborhood $N(v)\cap S_t$~in~$S_t$.
  \item Let $R_t \coloneqq\setof{ab\in E(G)}{a,b\in L_t}$ be the set of all edges in $G$ with both endpoints in $L_t$.
\end{itemize}
We verify that $(\cal P_1,R_1),\ldots,(\cal P_n,R_n)$ is a merge sequence of $G$. Clearly, $\cal P_1$ is the partition into singletons, 
and $\cal P_n$ has one part.
Observe that for $t\in[n]$ and any two parts $A,B$ of $\cal P_t$,
either $A,B\subset L_t$, and therefore  $E(G)\cap AB\subset R_t$,
or, otherwise, at least one of $A,B$ is contained in $S_t$, and then $A$ and $B$ are homogeneous,
as both $A$ and  $B$ are atomic types over $S_t$. 
In any case, $AB-R_t\subset E(G)$ or 
$AB-R_t\subset {V\choose 2}-E(G)$, so the conditions of a merge sequence are satisfied.


Fix $t\in[n]$. Fix two vertices $v,w\in L_t$, such that 
there is a path of length at most $r$ in $(V,R_t)$ from $v$ to $w$.
Then there is a path $\pi$ of length at most \(r\) from $v$ to $w$ in $G$
contained in $L_t$.
In particular, for every vertex $u\in N_G(w)\cap S_t$ 
there is a path from $v$ to $u$ of length at most $r+1$ in \(G\) (namely the path $\pi$ followed by the edge $wu$) such that every inner vertex on that path is in $L_t$.
In other words, $N_G(w)\cap S_t \subseteq \wreach_{r+1}(G,\le,v)$.
 As $|\wreach_{r+1}(G,\le,v)|\le k$ by assumption, and 
 the atomic type of $w$ over $S_t$ is uniquely 
determined by $N(w)\cap S_t$, it follows that 
there are at most $2^k$ parts of $\cal P_t$ 
that are reachable by a path of length at most $r$ from $v$.
Thus, the radius-$r$ width of $(\cal P_t,R_t)$ is at most
$2^k$.



To bound the radius-$r$ width of $(\cal P_{t-1},R_t)$, for $t>1$, notice that every part in $\cal P_{t}$ 
is a union of at most three parts in $\cal P_{t-1}$.
 Namely, if $v$ is the unique vertex with $v\in S_{t-1}-S_{t}$, then the atomic type $A$ of a vertex $u$ over $S_{t}$ 
 is uniquely determined by the atomic type of $u$ over $S_{t-1}$ and the atomic type of $u$ over $\set v$. 
 There are three 
 possible atomic types over $\set v$, corresponding 
 to
 whether $u=v$, or $(u\neq v)\land E(u,v)$, or 
$(u\neq v)\land \neg E(u,v)$. This proves that $A$ is a union of at most three parts in $\cal P_{t-1}$.

It follows that the radius-$r$ width of $(\cal P_{t-1},R_t)$ is at most $3\cdot 2^k$, and we have constructed a merge sequence of $G$ 
of radius-$r$ width at most $3\cdot 2^k$.
\end{proof}


% Namely, let $v$ be the unique vertex with $v\in U_{i}-U_{i+1}$.
% Let $\cal P_{i}'$ be the partition obtained from $\cal P_i$ by merging 
% the part $\set{v}\cal P_i$ with the part $A:=\setof{w\in V(G)-U_i}{N(w)\cap U_i=N(v)\cap U_i}\cal P_i$
% (if $A=\emptyset$, we don't do any merge). 
% Now, the partition $\cal P_{i+1}$ is obtained from the partition $\cal P_{i}'$
% by merging some pairs of parts, in pairs.

% We need to merge pairs of distinct parts $A,A'\in\cal P_i'$ such that
% there is a set $B\subset U_{i+1}$ with:
% \begin{align}\label{eq:AA'B}
%   N(A)\cap U_i=B\qquad\text{and}\qquad N(A')\cap U_i.
% \end{align}
% Note that for each set $B$ as above, we have 
% that  $B\subset \wreach^{2}_{G,\le}(v)$.
% Indeed, as $A,A'$ are distinct, it must be the case 
% that either one of $A,A'$ is equal to $\set{v}$ (and then even $B\subset\wreach^{1}_{G,\le}(v)$),
% or one of $A,A'$ contains a neighbor $a$ of $v$.
% In the latter case, for any $b\in B$, the path $v-a-b$ 
% witnesses that $b\in \wreach^{2}_{G,\le}(v)$.
% Furthermore, the merged pair $(A,A')$ is uniquely determined by $B$ and by $v$,
% as $A$ and $A'$ are the unique atomic types over $U_i$ satisfying~\eqref{eq:AA'B}.
% In particular, the number of such merged pairs is at most $2^k$.

% As the radius-$r$ width of $(\cal P_i,R_i)$ is at most $2^k$,
% and $|\cal P_{i}\le |\cal P_{i+1}|+2^k+1$,
% it follows easily that 
% we can obtain a merge sequence of radius-$r$ width at most 
% $2^k+2^k+1=2^{k+1}+1$.






% A \emph{merge sequence} with vertex set $V$
% is a finite sequence of steps of the following form. Start with the partition $\cal P$ of $V$ into singletons, and all pairs in ${V\choose 2}$ marked as \emph{unresolved}. At each step, either 
% merge two parts of the partition $\cal P$ into one, 
% or \emph{resolve} two (possibly equal) parts $A,B$ of $\cal P$ \emph{positively} (resp. \emph{negatively}),
% by proclaiming each unresolved pair $\set{a,b}$ with $a\in A$ and $b\in B$, $a\neq b$, an \emph{edge} 
% (resp. a  \emph{non-edge});
% the pair $\set{a,b}$ is from now on considered \emph{resolved}.
% After the last step,  all pairs in $V\choose 2$ must be resolved. We have thus constructed a graph with vertex set $V$, and the  proclaimed edges. At every time step $t$ we have a partition $\cal P_t$ of $V$ and set $R_t\subset {V\choose 2}$ of resolved 
% pairs at time $t$, where the partitions $\cal P_t$ become coarser and coarser, while the sets $R_t$
% become larger and larger as the time progresses. 




% It can be easily observed that 
% it is unncessary to resolve a pair $A,B$,
% unless either $A$ or $B$ (or both) will participate in 
% the next merge, and $A$ 

% Equivalently, a merge sequence of a graph $G$ is a sequence
% $$(\cal P_1,R_1),\ldots,(\cal P_m,R_m),$$
% where  $\cal P_i$ is a partition of $V(G)$ and $R_i\subset {V(G)\choose 2}$, 
% with $R_1=\emptyset$ and $R_m={V(G)\choose 2}$,
% and such that for $i=2,\ldots,m$, either the partition $\cal P_i$ is obtained from $\cal P_{i-1}$ by merging some two parts into one and $R_i=R_{i-1}$, 
% or otherwise, $\cal P_i=\cal P_{i-1}$ and $R_i=R_{i-1}\cup \setof{\set{a,b}}{a\in A,b\in B,a\neq b}$ for some
% $A,B\in \cal P$, 
% with $R_i-R_{i-1}$ 
%  either contained in $E(G)$,
% or disjoint from $E(G)$.







% Fix a parameter $r\in\N\cup\set{\infty}$.
% The \emph{radius-$r$ width} of a merge sequence is 
% the maximum, over 
% all time steps $t$ and vertices $v\in V(G)$,
% of the number of parts $P\in\cal P_t$ that are \emph{$r$-reachable from $v$}, in the sense 
% that there is a path from $v$ to some vertex  $w\in P$ of length ${\le}r$ in the 
% graph $(V,R_t)$ (so the path can follow both edges and non-edges of the current step).
% The  \emph{radius-$r$ merge-width} of a graph $G$, denoted $\mw_r(G)$, is 
%   the smallest radius-$r$ width of a merge sequence of $G$.
%  Finally, a graph class $\CC$ has \emph{bounded merge-width}
% if for every (finite) $r\in\N$ there is a bound $c_r\in\N$
% such that $\mw(G)\le c_r$ for all $G\in\CC$.

 

% It is important to note that at any step and part $P\in\cal P_t$,
% even though for each fixed vertex $v\in P$ the number of $r$-reachable parts from $v$ 
% is bounded by $k$, the total number of parts that are reachable from some vertex $v\in P$
% is potentially unbounded.





\subsection{Merge-width and flip-width}\label{sec:fw}
We study the relationship between merge-width and flip-width, whose definition we now recall.

For a graph $G$ and 
two sets $A,B\subset V(G)$,
\emph{flipping} the pair $(A,B)$ in $G$ 
results in the graph $G'$ 
with edges $E(G')=E(G)\triangle AB$, where \(\triangle\) denotes the symmetric difference.
Given a partition $\cal P$ of $V(G)$, 
a \emph{$\cal P$-flip} of $G$ is a graph $G'$
obtained from $G$ by flipping arbitrary pairs $A,B\in\cal P$ (possibly with $A=B$).
Thus, a graph $G'$ is a $\cal P$-flip of $G$ if and only if for every pair $A,B\in\cal P$, 
either $E(G')\cap AB=E(G)\cap AB$ or $E(G')\cap AB=AB-E(G)$.
A \(k\)-flip is a \(\cal P\)-flip with \(|\cal P| \le k\).

We recall the definition of flip-width from \cite{flip-width}.
Fix $k,r\in\N$.
The \emph{flip-width game} of radius $r$ and width $k$ is played on a graph $G$ between two players, flipper and runner.
Initially, runner picks a vertex $v_0\in V(G)$.
In round $t=1,2,\ldots$,
    flipper announces a $k$-flip $G_t$ of $G$;
    then runner picks $v_t\in B^r_{G_{t-1}}(v_{t-1})$ -- that is, they traverse 
    a path of length at most $r$ in the \emph{previous} graph $G_{t-1}$.
    Flipper wins if $v_t$ is isolated in $G_t$.
The \emph{radius-$r$ flip-width} of $G$, denoted $\fw_r(G)$,
is the least number $k$ such that flipper has a winning strategy for the flip-width game of width $k$ and radius $r$ on $G$.
We prove:
\begin{theorem}\label{thm:fw-cases}
  Every class of bounded merge-width has bounded flip-width.
  More precisely, for all $r\in\N$ and every graph $G$,
  $$\fw_r(G)\le 4^{\mw_{2r-1}(G)}.$$
\end{theorem}

\Cref{thm:fw-cases} allows to deduce several results about classes of bounded merge-width
from the corresponding results for classes of bounded flip-width.
For instance, we obtain the two corollaries stated in the introduction:
\corbe*
The forward implication in \Cref{cor:be} follows from \Cref{thm:be} and the trivial fact that each class of bounded expansion excludes some biclique as a subgraph. The backwards direction follows from \Cref{thm:fw-cases} and from the analogous statement for classes of bounded flip-width, proved in \cite[Thm. VI.3]{flip-width}.


\cortww*
In \Cref{cor:tww-mw}, we use the notion of merge-width for binary structures, specifically for \emph{ordered graphs} --  graphs equipped with a total order on the vertex set. This is defined in \Cref{sec:computing}.  \Cref{thm:fw-cases} applies to classes binary structures, allowing to derive \Cref{cor:tww-mw} by a similar reasoning as for \Cref{cor:be} above:
The forward direction uses \Cref{thm:tww}and the fact that every graph can be extended with a total order, without increasing its twin-width \cite{tww4}, and the backwards direction uses \Cref{thm:fw-cases} and the analogue of \Cref{cor:tww-mw} proved in \cite[Thm. VII.3]{flip-width}.
For simplicity, below we present the proof of \Cref{thm:fw-cases} only for graph classes. The more general case is analogous, 
but the requires notion of flips and of flip-width for binary structures as defined in \cite[Sec. V.B]{flip-width}.

We state another corollary that follows from \Cref{thm:deg} and an analogous result from \cite[Thm. VI.1]{flip-width}.
\begin{corollary}\label{cor:deg}
  Let $\CC$ be a graph class. Then $\CC$ has bounded degeneracy if and only if $\mw_1(\CC)<\infty$ and $\CC$ excludes some biclique $K_{t,t}$ as a subgraph.
\end{corollary}




\medskip
To prove \Cref{thm:fw-cases},
we start with a simple observation that allows us to rephrase the main condition from the definition 
of a merge sequence in terms of flips, as follows.


\begin{lemma}\label{lem:flip}
    Let $G=(V,E)$ be a graph, let $\cal P$ be a partition of $V$, and let $R\subset {V\choose 2}$.
    The following conditions are equivalent:
    \begin{enumerate}
        \item for all $A,B\in\cal P$, either $AB-R\subset E$, or $AB-R\subset AB - E$,
        \item  there is a $\cal P$-flip $G'$ of $G$ with $E(G')\subset R$.
    \end{enumerate}
    \end{lemma}
    \begin{proof}
        (1$\rightarrow$2).
        The $\cal P$-flip $G'$
        is obtained from $G$ by flipping each pair $A,B\in\cal P$ such that $AB-R\subset E$. As a result $AB-R\subset AB-E(G')$, 
        and therefore $E(G')\cap AB\subset R$ for each $A,B\in\cal P_t$.
        Altogether, $E(G')\subset R$.
    
        (2$\rightarrow$1). Fix $A,B\in\cal P$. We have that
        $E(G')\cap AB\subset R\cap AB$,
        so $AB-R\subset AB-E(G')$.
        As $G'$ is a $\cal P$-flip of $G$,
        we have that either $E(G')\cap AB=E(G)\cap AB$ or $E(G')\cap AB=AB-E(G)$.
        Altogether, either $AB-R\subset E(G)\cap AB$, or  $AB-R\subset AB-E(G)$.
    \end{proof}
    
Using \Cref{lem:flip}, we may therefore redefine merge-width in terms of flips, as follows.


% Fix a graph $G$.
% A \emph{monotone flip-width strategy} for $G$
% is a sequence of steps of the following form.
% Start with $E:=E(G)$, $R:={V\choose 2}$, 
% and $\cal P:=\set{V(G)}$. 
% We call the set $R$ the \emph{restraint}.
% In each step, one of three operations is performed, under the condition that 
% the inclusion $E\subset R$ is maintained:
% \begin{itemize}
%     \item refine the partition $\cal P$ of $V(G)$, thus replacing $\cal P$ by its refinement;
%     \item flip $E$ with respect to some pair $A,B$ of sets in $\cal P$, thus replacing $E$ with $E\triangle AB$;
%     \item remove some set of pairs $\set{a,b}$ from $R$, thus replacing $R$ by some subset of $R-E$.    
% \end{itemize}
% The sequence terminates when $R=\emptyset$.
% In particular, as $E\subset R$, in the end we have $E=\emptyset$.

% Fix $r\in\N\cup\set{\infty}$.
% The \emph{radius-$r$ width} of the strategy 
% is the maximum, over all time steps of the radius-$r$ width of $(\cal P,R)$.
% Denote by $\mfw_r(G)$ the minimum radius-$r$ width of any monotone flip-width strategy for $G$.

\begin{definition}Let $G$ be a graph.
A \emph{restrained flip sequence} for $G$ is a sequence
\begin{align}\label{eq:mono-fw}
    (\cal P_1,R_1,G_1),\ldots,(\cal P_m,R_m,G_m)    
\end{align}
such that:
\begin{itemize}
    \item $\cal P_1,\ldots,\cal P_m$ is a refining sequence of partitions, starting with the  partition $\cal P_1$ of $V(G)$ with one part, and ending with the partition $\cal P_m$ into singletons,
    \item the graph $G_t$ is a $\cal P_{t}$-flip of $G$, for $t\in [m]$,
    \item ${V\choose 2}= R_1\supseteq \cdots \supseteq R_m=\emptyset$, and
    \item  $E(G_t)\subset R_t$ for $t\in[m]$.
\end{itemize}
The radius-$r$ width of the restrained flip sequence is 
the maximum, over $t\in[m-1]$,
of the radius-$r$ width of $(\cal P_{t+1},R_{t})$.
\end{definition}

Note that in the reformulation above,
the partitions $\cal P_1,\ldots,\cal P_m$ are becoming finer with each step, instead of becoming coarser as in merge sequences. In each step of a restrained flip sequence, we provide a $\cal P_t$-flip $G_t$.
The sequence $R_1\supseteq \cdots\supseteq R_m$
is a descending sequence of subsets of $V\choose 2$ with $R_m=\emptyset$. The set $R_t$ is called the \emph{restraint} at time $t$, as we require that $E(G_s)\subset R_t$ for all $s\in[m]$ with $s\ge t$, thus restraining all the future graphs $G_t,G_{t+1},\ldots,G_m$.  Note that  $E(G_m)\subset R_m=\emptyset$, so $G_m$ is edgeless.

As $G_t$ is a $\cal P_t$-flip of $G$, 
$G_{t-1}$ is a $\cal P_{t-1}$-flip of $G$, and $\cal P_{t-1}$ is coarser than $\cal P_t$,
we may equivalently require that $G_t$ as a $\cal P_{t}$-flip of $G_{t-1}$, rather than of $G$.
Therefore, in the sequence $G_1,G_2,G_3,\ldots$ 
each subsequent graph is a flip of the previous graph.
This is similar to the setting of flipper games considered in \cite{flippers}.

% Note that for the second operation to be applicable, to ensure $E\subset R$, it must be the case that $\set{a,b}\in R$ for all $a\in A$ and $b\in B$.
% Similarly, for the third operation to be applicable, it must be the case that  $\set{a,b}\in R-E$.



\begin{lemma}\label{lem:mw-flip seq}
    For every graph $G$,
    there is a correspondence between merge sequences and restrained flip sequences for $G$,
    which preserves the length, and the radius-$r$ width of the sequences, for each $r\in\N\cup\set{\infty}$.
\end{lemma}




\begin{proof}
    Fix a merge sequence 
    $(\cal P_1,R_1),\ldots,(\cal P_m,R_m)$
    of $G$.
 By \Cref{lem:flip} (1$\rightarrow$2),
 for each $t\in[m]$,
    there is a  $\cal P_t$-flip $G_t$ of $G$ such that 
$E(G_t)\subset R_t$.
    The sequence 
    $(\cal P_m,R_m,G_m),\ldots,(\cal P_1,R_1,G_1)$
    is a restrained flip sequence for $G$.
    
    Conversely,
    given a restrained flip sequence for $G$ of the form $(\cal P_1,R_1,G_1),\ldots,(\cal P_m,R_m,G_m)$,
    by \Cref{lem:flip} (2$\rightarrow$1) we conclude that 
    the sequence $(\cal P_m,R_m),\ldots,(\cal P_1,R_1)$ is a merge sequence of $G$.
\end{proof}

% \begin{remark}\label{rem:isolations}\sz{rewrite, reorganize this stuff}
% Define a \emph{restrained isolation sequence} for $G$ 
% as a sequence 
% $$G_0,G_1,G_2,\ldots,G_n,$$
% where 
% $G_0=G$ and for all $t=1,\ldots,n$, 
% $G_{t}$ is obtained from the graph $G_{t-1}$ by \emph{isolating} 
% some vertex $v_{t}\in V(G_t)-\set{v_1,\ldots,v_{t-1}}$ (that is, removing all edges incident to $v_t$). 

% Then $G_t$ is a $\cal P_t$-flip of $G$,
% where $\cal P_t$ is the partition of $V(G)$ into atomic types over the set $\set{v_1,\ldots,v_{t}}$ (see proof of \Cref{lem:wcol}).
% Therefore, 
% $$(\set{V},\mbox{$V\choose 2$},G),\ (\cal P_1,E(G_1),G_1),\ldots,(\cal P_n,E(G_n),G_n)$$
% forms a restrained flip sequence.
% The radius-$r$ width of this sequence is functionally equivalent to the medium $(r+1)$-reachability number of $G$,
% as defined in \Cref{sec:sparsity},
% where the total order on the vertices corresponds to the enumeration $v_1<\ldots<v_n$.

% Therefore, the  medium $(r+1)$-reachability number of $G$ 
% (related to generalized coloring numbers 
% via \Cref{lem:cols})
% can be viewed as defining a special form of restrained flip sequences, and therefore, of merge-sequences. In particular, by \Cref{fact:be}, for the limit value $r=\infty$, restrained isolation sequences recover (up to functional equivalence), the treewidth parameter.
% \end{remark}
% \jan{this has changed when removing medium reachability.}


% The \emph{radius-$r$ width} of the strategy is defined similarly, but instead, counts the number 
% of sets in the partition of $V(G)$ generated by $\cal F$ (that is, the atoms of the boolean subalgebra of $2^{V(G)}$ generated by $\cal F$, or the partition corresponding to the equivalence relation on $V(G)$,
% where two vertices are equivalent if and only if they belong to the same sets in $\cal F$).




% Say that a graph $G'=(V,E')$ is \emph{$R$-locally} a $k$-flip of $G=(V,E)$  
% if for every $v\in V$ there is a $k$-flip $G''$ of $G$ such that $B^r_G'(w)=B^r_{G'}(w)$ holds for all $w\in B^r_{R}(v)$.

The next lemma shows that radius-$r$ flip-width is bounded in terms of radius-$(2r-1)$ merge-width, thus proving \Cref{thm:fw-cases}.
% It also suggests that restrained flip sequences may be viewed as monotone winning strategies for the pursuer in the flip-width game.\sz{reference}




\begin{lemma}\label{lem:fw}
    Fix $r,s\in\N$. 
    For every graph $G$, if $\mw_{2r-1}(G)\le s$ then  $\fw_r(G)\le 4^s$. 
    % Moreover, 
    % there is a sequence $R_1\supseteq \cdots \supseteq R_n=\emptyset$ such that 
    % flipper has a strategy of width $4^s$ and radius $r$ which wins in $|V(G)|$ steps on $G$,
    % and such that $B^r_{G_t}(v_t)\subset B^r_{R_t}(v_t)$ holds for all $t\in [n]$.
\end{lemma}

\begin{proof}
    Let $G$ be a graph with $\mw_{2r-1}(G)\le s$. Then $G$ has a merge sequence of radius-$(2r-1)$ merge-width at most $s$.
    By \Cref{lem:mw-flip seq}, $G$ has a 
 restrained flip sequence of radius-$(2r-1)$ width at most \(s\) of the form
$$(\cal P_1,R_1,G_1),\ldots,(\cal P_n,R_n,G_n).$$

The following lemma will allow us to construct a strategy for flipper in the flip-width game.

\begin{lemma}\label{lem:strategy}
    Fix $t\in [2,n]$, and 
    let $s$ be the radius-$(2r-1)$ width of $(\cal P_t,R_{t-1})$.
    For every $v\in V(G)$ there is a $4^s$-flip $G_t'$ of $G$ such that:
$$B^r_{G_t'}(w)\subset B^r_{G_t}(w)\quad\text{for every $w\in B^r_{G_{t-1}}(v)$.}$$
\end{lemma}
Before proving \Cref{lem:strategy}, we first show how it yields a winning strategy
for flipper in the flip-width game of radius $r$
    and width $4^s$.
Flipper's strategy will be to announce graphs $G_1',G_2',\ldots$, ensuring  that following invariant holds after round $t\ge 1$ of the game:
    \begin{align}\label{eq:fw-invariant}
    B^r_{G_{t}'}(v_t)\subset B^r_{G_t}(v_t).    
    \end{align}
    
    In the first round, when $t=1$, flipper announces $G_1'=G_1$ and runner picks a vertex 
    $v_1$,
    and the invariant holds trivially.
    Suppose the invariant is satisfied after round $t-1$ of the game,
    that is, 
    $$B^r_{G_{t-1}'}(v_{t-1})\subset B^r_{G_{t-1}}(v_{t-1}).$$
    Now flipper announces the $4^s$-flip $G_t'$ of $G$
    given by \Cref{lem:strategy} for $v \coloneqq v_{t-1}$.
    Next, runner picks a vertex $v_t\in B^r_{G_{t-1}'}(v_{t-1})$.
In particular, $v_t\in B^r_{G_{t-1}}(v_{t-1})$, so 
 \eqref{eq:fw-invariant} holds by \Cref{lem:strategy}, and the invariant is fulfilled.

    Playing according to this strategy, flipper wins within $n$ rounds,
    since $E(G_n)=\emptyset$, and therefore $v_n$ is isolated in $G_n'$
    by \eqref{eq:fw-invariant}. This proves that $\fw_r(G)\le 4^s$, and thus \Cref{lem:fw}.
\end{proof}

\begin{remark}\label{rem:duration}
    Observe that the number of rounds needed by flipper to win in the flip-width game of radius $r$ and width $4^s$ can be bounded by $|V(G)|$.
    Namely, the proof of \Cref{lem:fw} above shows that 
    the number of rounds is (at most) equal to the length $n$ of a restrained flip sequence 
     of radius-$(2r-1)$ width at most \(s\) for $G$.
     By \Cref{lem:mw-flip seq}, this corresponds to the length of a radius-$(2r-1)$ merge sequence of width $s$. Any merge sequence of $G$ can be converted into one of length $n=|V(G)|$, while preserving its radius-$r$ width for each $r\in\N$, since if for two consecutive pairs $(\cal P_i, R_i),(\cal P_{i+1},R_{i+1})$ in a merge sequence we have $\cal P_i=\cal P_{i+1}$, then $(\cal P_{i+1},R_{i+1})$ can be dropped from the merge sequence.
\end{remark}

We now prove \Cref{lem:strategy}. 



    \begin{proof}[Proof of \Cref{lem:strategy}]
        Fix $v\in V(G)$. Let
 $\cal N\subset \cal P_t$ consist of all parts $A\in\cal P_t$ that can be reached by a path of length at most $2r-1$ by from $v$ in the graph $(V,R_{t-1})$.
In particular, $|\cal N|\le s$.

Let $G_t'=G'=(V,E)$ be the graph 
such that for any two parts $A,B\in\cal P_t$:
$$E'\cap AB=
\begin{cases}
    E(G)\cap AB&\text{if $A,B\notin\cal N$}\\
    E(G_t)\cap AB&\text{otherwise}.
\end{cases}$$

\begin{claim}\label{cl:k-flip}
    $G'$ is a $4^s$-flip of $G$.
\end{claim}
\begin{claimproof}
    As $G_t$ is a $\cal P_t$-flip of $G$, so is $G'$.
    Therefore, for any two parts $A,B\in\cal P_t$, 
    the set $E'\cap AB$ is either equal to $E\cap AB$,
    or to $AB-E$.

    For every part $A\in\cal N$, let $U(A)$ be the union of all parts 
    $B\in\cal P_t$ such that $E'\cap AB= AB-E$.
    Then for every $B$ with $B\subset U(A)$ or $B\subset V- U(A)$,
    we have that
    $E'\cap AB=E\cap AB$ or 
     $E'\cap AB=AB-E$.

    Consider the set family:
    $$\cal F\coloneqq \cal N \cup \setof{U(A)}{A\in \cal N}.$$ Then $|\cal F|\le 2|\cal N|\le 2s$.
    
    Let $\cal Q$ be the partition 
    of $V$ such that 
    two vertices $a,b$ are in the same part of $\cal Q$ if and only if $a$ and $b$ belong to the same sets in $\cal F$.
    Then $|\cal Q|\le 2^{|\cal F|}\le 4^{s}$.
    Moreover, for every $A,B\in \cal Q$,
    either $E'\cap AB=E\cap AB$,
    or $E'\cap AB=AB-E$.
    It follows that 
    $G'$ is a $4^s$-flip of $G$.
\end{claimproof}



\begin{claim}\label{cl:balls}
    For all $w\in B^r_{G_{t-1}}(v)$
 and $i=0,\ldots,r$ we have:
    $$B^i_{G'}(w)\subset B^i_{G_t}(w).$$
\end{claim}

\begin{claimproof}
    We induct on $i$. For $i=0$ the statement is trivial.
In the inductive step, fix $1\le i\le r$ and $u\in B^i_{G'}(w)$ with $u\neq w$.
Then there is some  $a\in B^{i-1}_{G'}(w)$ with $ua\in E(G')$.
By inductive assumption, $a\in B^{i-1}_{G_t}(w)$.
From  $i\le r$ and  $E(G_t)\subset R_{t-1}$
we get that $a\in B^{r-1}_{R_{t-1}}(w)$,
which together with $w\in B^{r}_{G_{t-1}}(v)\subset B^{r}_{R_{t-1}}(v)$ implies $a\in B^{2r-1}_{R_{t-1}}(v)$. In particular, $a\in \bigcup \cal N$.
By definition of $E'$, this implies $ua\in E(G_t)$. Together with 
$a\in B^{i-1}_{G_t}(w)$, this implies that $u\in B^{i}_{G_t}(w)$, as required.
\end{claimproof}
\Cref{cl:k-flip} and \Cref{cl:balls} together prove \Cref{lem:strategy}.
\end{proof}




\subsection{Almost bounded merge-width}\label{sec:abmw}
A graph class $\CC$ has \emph{almost bounded merge-width}
if for every $r\in\N$ and $\eps>0$,
we have that $\mw_r(G)\le O_{\CC,r,\eps}(|V(G)|^\eps)$, for all $n\in\N$ and all $n$-vertex graphs $G\in\CC$.

Clearly, classes of bounded merge-width have almost bounded merge-width.
The following result states that all nowhere dense classes have almost bounded merge-width. We will prove this later below, by inspecting the proof of \Cref{thm:be}.

\thmnwd*

Classes of \emph{almost bounded flip-width} are defined analogously as classes of almost bounded merge-width, with $\mw_r(G)$ replaced with $\fw_r(G)$.
The following is the main result of this section.

\thmabmw*

As every hereditary class of almost bounded flip-width is monadically dependent \cite[Thm. 2.12]{flip-breakability}, we obtain the following.
\begin{corollary}\label{cor:abmw-mNIP}
    Every hereditary class of almost bounded merge-width is monadically dependent.
\end{corollary}

\Cref{thm:nwd} and \Cref{cor:abmw-mNIP} together justify our \Cref{conj:almostboundedmergewidth}
that almost bounded merge-width coincides with monadic dependence on hereditary classes.


% We give a short proof of \Cref{fact:nbd}, using \Cref{fact:nd-wcol}.
% \begin{proof}
%     Fix $G\in\CC$, and let $n=|V(G)$. Let $\le$ be a total order on $V(G)$ 
%     such that every vertex $v\in V(G)$ weakly $2$-reaches at most $n^\eps$ vertices in $V(G)$, with respect to the order $\le$. 
%     In particular, for every $v\in V(G)$, the set $N_G^<(v):=\setof{w\in N_G(v)}{w\le v}$
%     has size at most $n^\eps$.


% Let $S\subset V(G)$ be arbitrary.
% We cover $V(G)$ with three sets:
%     \begin{itemize}
%         \item $S$
%         \item $T:=\bigcup_{s\in S}N_G^<(s)$,
%         \item $U:=V(G)-(S\cup T)$.
%     \end{itemize}
%     We have that $|S|+|T|\le |S|+|S|\cdot n^\eps\le |S|(1+n^\eps)$. 
% We bound the 
% size of the family
% $$\setof{N_G(v)\cap S}{v\in U}.$$ 

% Note that for $v\in U$ we have that $N_G(v)\cap S=N_G^<(v)$,
% since $v\notin T$, so $v$ has no neighbors $s\in S$ with $s>v$.

% For $v\in U$ with $N_G(v)\cap S\neq\emptyset$, let $m_S(v):=\max N_G(v)=\max N_G^<(v)$.
% Then $m_S(v)<v$ and $m_S(v)\in S$. Moreover,
% $N_G^<(v)\subset \wreach_2^{G,\le}(m_S(v))$.


% \end{proof}


\subsection{Preliminaries}
Before proving \Cref{thm:nwd} and \Cref{thm:abmw}, 
we first recall some basic notions.

\medskip
A \emph{set system} is a pair $(X,\cal F)$ with $\cal F\subset 2^X$.
Its \emph{VC-dimension} is the maximal size of a subset $Y\subset X$ such that 
$\setof{Y\cap F}{F\in\cal F}=2^Y$.
We recall the fundamental Sauer-Shelah-Perles lemma~\cite{sauer,shelah-sauer-lemma}.

\begin{lemma}[Sauer-Shelah-Perles lemma]\label{lem:sauer-shelah-perles}
  Let $(X,\cal F)$ be a set system of VC-dimension~$d$.
  Then $|\cal F|\le O(|X|^d)$.
\end{lemma}

The \emph{VC-dimension} of a graph $G$, denoted $\VCdim(G)$, 
is defined as the VC-dimension of the set system $\bigl(V(G),\setof{N(v)}{v\in V(G)}\bigr)$.
More explicitly, $\VCdim(G)$ is the maximal size of a subset $X\subset V(G)$
such that $\setof{N(v)\cap X}{v\in V(G)}=2^X$.

Define the \emph{neighborhood complexity} of a graph $G$, as the function $\pi_G\from\N\to\N$ defined as
 $$\pi_G(s)\coloneqq\max\Bigl\{\bigl|\{N_G(v)\cap S\,:\,v\in V(G)\}\bigr|\,:\, S\subset V(G), |S|\le s\Bigr\} \quad\text{for all }s \in \N.$$
% (Equivalently, $\pi_G(s)$ is the maximal number of atomic types over a set $S$ of size at most $s$ in $G$. This is essentially equal to the \emph{neighborhood complexity}, up to an additive factor of $s$.) 
In particular, $\pi_G(s)\le 2^s$, for all graphs $G$ and all $s\in\N$.
From \Cref{lem:sauer-shelah-perles}, we get the following.
\begin{corollary}\label{cor:sauer-shelah}
    Let $G$ be a graph of VC-dimension $d$.
    Then $\pi_G(s)\le O(s^d)$, for all $s\in \N$.
\end{corollary}



\subsection{Proof of \Cref{thm:nwd}}


To prove \Cref{thm:nwd}, we use the following result.
\begin{fact}[\cite{zhu2009colouring,sparsity-book}]\label{fact:nd-wcol}
    Let $\CC$ be a nowhere dense graph class, and fix $r\in\N$ and $\eps>0$.
    Then for every graph $G\in\CC$
    $$\wcol_r(G)\le O_{\CC,r,\eps}(|V(G)|^{\eps}).$$
\end{fact}

% \begin{fact}[\cite{eickmeyer_et_al:LIPIcs.ICALP.2017.63}]\label{fact:nbd}
%     Let $\CC$ be a nowhere dense graph class and $\eps>0$.
%     Then for every graph $G\in\CC$ we have:
%     $$\pi_G(s)\le O_{\CC,\eps}(|V(G)|^\eps),\qquad\text{for all $s\in\N$}.$$
% \end{fact}
\begin{lemma}\label{lem:nd-VC}
    Let $\CC$ be a nowhere dense graph class. 
    Then there is some $d\in \N$ such that all graphs $G\in \CC$ have VC-dimension less than $d$.
\end{lemma}
\begin{proof}
    As $\CC$ is nowhere dense, there is some $d\in\N$ such that 
    no graph $G\in \CC$ contains the $1$-subdivision of the clique $K_d$, as a subgraph.
    It follows that every graph $G\in \CC$ has VC-dimension less than $d$.
\end{proof}


\begin{lemma}\label{lem:wcol-poly}
    For every $r\in\N\cup\set{\infty}$ and graph $G$, we have
    $$\mw_r(G)\le 3\cdot\pi_G(\wcol_{r+1}(G)).$$
  \end{lemma}
  \begin{proof}[Proof sketch]
    We follow the proof of \Cref{lem:wcol}, and use the same notation. 
    Let \(k\coloneqq\wcol_{r+1}(G)\).
    Fix $t\in [n]$, and let $L_t,S_t,\cal P_t,R_t,\le$ be as defined in the proof.
    Fix $v\in L_t$.
    It was observed that for every $w\in V(G)$ which is reachable from $v$ by a path of length $r$
in the graph $(V,R_t)$, 
 the atomic type of $w$ over $S_t$ is uniquely 
determined by $N_G(w)\cap S_t \subseteq \wreach_{r+1}(G,\le,v)$,
and moreover, $|\wreach_{r+1}(G,\le,v)|\le k$.
Since $$\Big|\setof{N_G(w)\cap \wreach_{r+1}(G,\le,v)}{w\in V(G)}\Big|\le  \pi_G(k),$$
we conclude that are at most $\pi_G(k)$ parts of $\cal P_t$ 
that are reachable by a path of length at most $r+1$ from $v$.
Thus, the radius-$r$ width of $(\cal P_t,R_t)$ is at most
$\pi_G(k)$. The rest of the argument remains unchanged, yielding the final bound $\mw_r(G)\le 3\cdot \pi_G(k)\le 3\cdot \pi_G(\wcol_{r+1}(G))$.
  \end{proof}
  
  \Cref{thm:nwd} now follows by combining the previous insights.
  \begin{proof}[Proof of \Cref{thm:nwd}]
    By \Cref{lem:nd-VC}, there is some $d\in\N$ such that every $G\in\CC$ has VC-dimension at most $d$.
    Fix $r\in\N$ and $\eps>0$.
Then, for every graph $G\in\CC$, by \Cref{lem:wcol-poly} and \Cref{fact:nd-wcol}, we have 
$$\mw_r(G)\le 3\cdot \pi_G\bigl(\wcol_{r+1}(G)\bigr)\le  3\cdot \pi_G\bigl(O_{\CC,r,\eps}(|V(G)|^\eps)\bigr)\le O_{\CC,r,\eps}\bigl(|V(G)|^{\eps\cdot d}\bigr).$$
Since $\eps>0$ is arbitrary, this implies that $\CC$ has almost bounded merge-width.
  \end{proof}

  \subsection{Proof of \Cref{thm:abmw}}


We start by proving that hereditary classes of almost bounded merge-width have bounded VC-dimension. 
In fact, we prove a stronger result, expressed using the following notion.

\newcommand{\ntn}{\mathrm{ntn}}

A pair $u,v$ of distinct vertices in a graph $G$ is a pair of \emph{$k$-near-twins} if $|N_G(u)\triangle N_G(v)|\le k$.
 For a bipartite graph $G=(X,Y,E)$, let the near-twin number $\ntn(G)$ denote 
 the smallest number $k$ such that 
for all $X'\subset X,Y'\subset Y$ with $|X'|+|Y'|>2$,
the bipartite subgraph $G[X',Y']$ induced by $X'$ and $Y'$  in $G$
has a pair of $k$-near-twins contained either in $X'$, or in $Y'$.

\begin{lemma}\label{lem:twins-bip}
    Let $G=(X,Y,E)$ be a bipartite graph with $\mw_1(G)\le k$
    and with $|X|+|Y|>2$.
    Then $G$ contains a pair of $2k$-near-twins
     contained in a single part of $G$. In particular, $\ntn(G)\le 2k$.
\end{lemma}
\begin{proof}
    Fix a construction sequence of $G$ of radius-$1$ width $k$. Consider the first moment when some part $A$ of the current partition $\cal P$ contains two vertices $a,b$ that belong to the same part of the bipartition $\set{X,Y}$ of $G$. Suppose $a,b\in X$. Then $N_G(a)\triangle N_G(b)\subset (N_R(a)\cup N_R(b))\cap Y$,
    where $N_R(\cdot)$ denotes the neighborhood in the graph $(V,R)$  of currently resolved pairs.
    Note that $|N_R(a)\cap Y|\le k$,
    since no two vertices of $Y$ are in a single part of the current partition $\cal P$, and $N_R(a)$ is contained in at most $k$ parts of $\cal P$. Similarly, $|N_R(b)\cap Y|\le k$.
    It follows that $|N_G(a)\triangle N_G(b)|\le 2k$.
\end{proof}

The following lemma is implicit in \cite{flip-width}
(it follows from 
\cite[Lem. 5.25]{flip-width-arxiv}).
\begin{lemma}\label{lem:ntn}
    If $\CC$ is a hereditary class of bipartite graphs such that $\ntn(G)\le o(|V(G)|)$ for $G \in \CC$, then $\VCdim(\CC)<\infty$.
\end{lemma}



% \begin{lemma}[{\cite[Lem. 5.25]{flip-width-arxiv}}]\label{lem:VC-no-twins}
%     Let $G$ be a graph with $\VCdim(G) \ge 2m$, for some $m$. Then there are two sets $X,Y\subset V(G)$ such that the bipartite graph $G[X, Y]$\sz{notation} contains no pair of $(2^{m-1} - 1)$-near-twins in either of the parts $X, Y$.
% \end{lemma}

% \begin{lemma}\label{lem:mw-vc}
%     Every graph $G$ with $\VCdim(G)\ge d$ contains an induced subgraph with $O(d)$ vertices and with $\mw_1(G)\ge d/8$.
% \end{lemma}
% \begin{proof}
%     Follows from \Cref{lem:twins-bip} and \Cref{lem:VC-no-twins}.\sz{relate $\mw_1(G[X,Y])$ with $\mw_1(G)$}
% \end{proof}

\begin{corollary}\label{lem:abmw-vc}
    If $\CC$ is a hereditary class of graphs such that $\mw_1(G)\le o(|V(G)|)$ for $G \in \CC$, then $\VCdim(\CC)<\infty$.
\end{corollary}
    \begin{proof}
For a graph $G$ define the bipartite graph $B(G)$, with two parts of size $V(G)$,
representing the binary relation $E(G)\subset V(G)\times V(G)$.
Then $\mw_1(B(G))\le O (\mw_1(G))$, as a construction sequence of $G$ can
be easily converted to a construction sequence for $B(G)$.
In particular, $\mw_1(B(G))\le o(|V(G)|)$ for $G\in \CC$.
Moreover, $\VCdim(G)=\VCdim(B(G))$.
The conclusion follows from \Cref{lem:twins-bip}
and \Cref{lem:ntn}, applied to the class $\setof{B(G)}{G\in\CC}$.
        % If $\VCdim(\CC) = \infty$ then for every $d$ there is a graph $G\in\CC$ with $\VCdim(G) \ge  d$, and by \Cref{lem:mw-vc} there is $H \in \CC$ with $O(d)$ vertices and $\mw_1(H) = \Omega(d)$. Since this holds for all $d \in\N$, it cannot be that $\mw_1(G) \le o(|V(G)|)$ for all $G \in \CC$.
    \end{proof}




    Recall that by the Sauer-Shelah-Perles Lemma (see \Cref{lem:sauer-shelah-perles}), 
    if $G$ is a graph of VC-dimension at most $d$,
    then $\pi_G(s)\le O(s^d)$ for all $s\in\N$. The following variant of \Cref{lem:fw}
     therefore gives -- for graphs of fixed VC-dimension -- a polynomial bound on the flip-width parameters, in terms of the merge-width parameters.

    
    \begin{lemma}\label{lem:fw-poly}Fix $r\ge 1$.
        For every graph $G$, if 
        $s=\mw_{2r-1}(G)$ then        
        $$\fw_r(G)\le s^2+\pi_G(s).$$
    \end{lemma}
\Cref{lem:fw-poly} strengthens 
 \Cref{lem:fw} by replacing the upper bound $4^s$ by $s^2+\pi_G(s)$.
    % However, \Cref{lem:fw-poly} gives a slightly better dependency on the radius, as the upper bound involves $\mw_{2r-1}(G)$, rather than $\mw_{2r}(G)$.
    Then \Cref{lem:fw-poly} follows from the lemma below, 
    exactly in the same way as \Cref{lem:fw} follows from \Cref{lem:strategy}.
    
    \begin{comment}
        
    
    \begin{lemma}\label{lem:strategy-poly}
        Fix $t\in [2,n]$, and 
        let $s$ be the radius-$2r$ width of $(\cal P_t,R_{t-1})$.
        For every $v\in V(G)$ there is a $\pi_G(s)$-flip $G_t'$ of $G$ such that:
    $$B^r_{G_t'}(w)\subset B^r_{G_t}(w)\quad\text{for every $w\in B^r_{G_{t-1}}(v)$.}$$
    \end{lemma}
    
    
    
    
    %     \begin{proof}Denote $s\coloneqq \mw_{2r}(G)$ and $V\coloneqq V(G)$.
    %         Fix $v\in V$. 
    %         For $i=0,\ldots,2r$, let 
    %         $B_i$ denote the set of vertices  that can be reached by a path of length at most $i$ by from $v$ in the graph $(V,R_{t-1})$.
    %         Let $\cal Q_i=\setof{A\cap B_{2r}}{A\in\cal P_t,A\cap B_{2r}\subset B_i}$.
    % In particular, $|\cal Q_{2r}|\le s$.
    
    
    % Observe that
    % there are no edges in $G_{t}$ between  $B_{2r-1}$ and $V-B_{2r}$,
    % since an edge  $uw\in E(G_t)$ with  $u\in B_{2r-1}$  implies $w\in B_{2r}$,
    % as $E(G_{t})\subset R_t\subset R_{t-1}$.
    % Since each part of $\cal Q_{2r-1}$ is contained in $B_{2r-1}$ and in some part of $\cal P_t$,
    %  and $G_t$ is a $\cal P_t$-flip of $G$, 
    % it follows that each part of $\cal Q_{2r-1}$ is homogeneous in $G$ towards 
    % every vertex $w\in V(G)-B_{2r}$.
    
    
    % Let $\cal R$ be the partition of $V(G)-B_{2r}$ which partitions vertices according to their neighborhood in $B_{2r-1}$, in $G$.
    % Denote $\cal Q\coloneqq \cal Q_{2r}\cup\cal R$. Then $\cal Q$ is a partition of $V$.
    
    % \begin{claim}\label{cl:small-q-poly}
    %     $|\cal Q|\le \pi_G(s).$
    %    \end{claim}
    %    \begin{claimproof}
    %        For every part $A\in \cal Q_{2r-1}$ pick a vertex $v_A\in B_{2r-1}\cap A$,
    %        and let $S=\setof{v_A}{A\in\cal Q_{2r-1}}$. In particular, $|S|\le s$.
    %        It follows from the above that, in the graph $G$, for every vertex $w\in V(G)-B_{2r}$,
    %        the neighborhood  of $w$ in $B_{2r-1}$ is uniquely determined by the neighborhood of $w$ in $S$; moreover, $V(G)-B_{2r}$ is disjoint from $S$.
    %        Therefore, $$|\cal R|\le  \pi_G(s)-s.$$
    %        Together with $|\cal Q_{2r}|\le s$, this yields the conclusion.
    %    \end{claimproof}
           
    
    % Consider the graph $G'$ with vertex set $V$,
    % such that for any two parts $A,B\in\cal Q$:
    % $$E(G')\cap AB=
    % \begin{cases}
    %     E(G)\cap AB&\text{if $A,B\notin \cal Q_{2r-1}$},\\
    %     E(G_t)\cap AB&\text{otherwise.}\\
    %     % \emptyset &\text{if $A\in \cal Q_{2r-2}$ and $B\notin  \cal Q_{2r-1}$}.
    % \end{cases}$$
    
    % \begin{claim}\label{cl:k-flip-poly}
    %     $G'$ is a $\cal Q$-flip of $G$.
    % \end{claim}
    % \begin{claimproof}
        
    %     It is enough to show that for all   $A,B\in \cal Q$, 
    %     \begin{align}
    %         E(G')\cap AB=E(G)\cap AB\qquad\textit{or}\qquad
    %         E(G')\cap AB= AB-E(G).\label{eq:flip-poly}    
    %     \end{align}
    %     We consider several cases.
    % \begin{itemize}
    %     \item If $A,B\notin \cal Q_{2r-1}$, this is clear by definition of $E(G')$.
    %     \item Suppose that $A\in\cal Q_{2r-1}$ and $B\in \cal Q_{2r}$.
    %     % Let $A',B'$ be the unique parts of $\cal P_t$ with $A\subset A'$ and $B\subset B'$.
    %     % Since  is a $\cal P_t$-flip of $G$,
    %     % we have that $E(G_t)\cap A'B'$ is either equal to $E(G)\cap A'B'$, or to $A'B'-E(G)$.
    %     Then \eqref{eq:flip-poly} follows, as $E(G')\cap AB=E(G_t)\cap AB$ by definition,
    %     and $G_t$ is a $\cal P_t$-flip of $G$, and each of the parts $A,B$ is contained in some part of $\cal P_t$.
    %     \item Suppose that $A\in\cal Q_{2r-1}$ and $B\notin\cal Q_{2r}$. Then 
    %     $A\subset B_{2r-1}$ and 
    %     $B\subset V(G)-B_{2r}$, and $B\in\cal R$. By the previous discussion, $E(G_t)\cap AB=\emptyset$. Moreover, each vertex of $b\in B$ is homogeneous in $G$ towards $A$,
    %     that is, each $b$ is either complete, or anti-complete towards $A$ in $G$.
    %     Since all vertices of $B$ have equal neighborhoods in $A$
    %     (by definition of $\cal R$, as $B\in\cal R$ and $A\subset B_{2r-1}$),
    %     it follows that $B$ is homogeneous in $G$ towards $A$.
    % Together with $E(G_t)\cap AB=\emptyset$, this yields \eqref{eq:flip-poly}.
    % \end{itemize}
    % Up to exchanging the roles of $A$ and $B$, this covers all the cases.
    % \end{claimproof}
    
    
    
    
    % \begin{claim}\label{cl:balls-poly}
    %     For all $w\in B^r_{G_{t-1}}(v)$
    %  and $i=0,\ldots,r$ we have:
    %     $$B^i_{G'}(w)\subset B^i_{G_t}(w).$$
    % \end{claim}
    
    % \begin{claimproof}
    %     We induct on $i$. For $i=0$ the statement is trivial.
    % In the inductive step, fix $1\le i\le r$ and $u\in B^i_{G'}(w)$ with $u\neq w$.
    % Then there is some  $a\in B^{i-1}_{G'}(w)$ with $ua\in E(G')$.
    % By inductive assumption, $a\in B^{i-1}_{G_t}(w)$.
    % From  $i\le r$ and  $E(G_t)\subset R_{t-1}$
    % we get that $a\in B^{r-1}_{R_{t-1}}(w)$,
    % which together with $w\in B^{r}_{G_{t-1}}(v)\subset B^{r}_{R_{t-1}}(v)$ implies $a\in B^{2r-1}_{R_{t-1}}(v)=B_{2r-1}$. In particular, $a\in \bigcup \cal Q_{2r-1}$\sz{error?}.
    % As $ua\in E(G')$, this implies $ua\in E(G_t)$, by definition of $E(G')$. Together with 
    % $a\in B^{i-1}_{G_t}(w)$, this implies that $u\in B^{i}_{G_t}(w)$, as required.
    % \end{claimproof}
    % \Cref{cl:k-flip-poly,cl:small-q-poly,cl:balls-poly} together prove \Cref{lem:strategy-poly}.
    % \end{proof}
% \begin{szin}
%     \textbf{This proof is new and probably flawed, fix below:}
%     \begin{proof}
%         Fix $v\in V(G)$. For $i=0,\ldots,2r$, let 
%  $\cal Q_i\subset \cal P_t$ consist of all parts $A\in\cal P_t$ that can be reached by a path of length at most $i$ by from $v$ in the graph $(V,R_{t-1})$.
% In particular, $|\cal Q_{2r}|\le s$.

% Let $G_t'=G'$ be the graph such that for any two parts $A,B\in\cal P_t$:
% $$E(G')\cap AB=
% \begin{cases}
%     E(G)\cap AB&\text{if $A,B\notin \cal Q_{2r-1}$},\\
%     E(G_t)\cap AB&\text{otherwise}.
% \end{cases}.
% $$
% Note that this graph is exactly the same as the graph defined in the proof of \Cref{lem:strategy}.
% % for $A,B\in \cal P_t$ and $a\in A,b\in B$,
% % $ab\in E(G')$ if and only if one of the following holds:
% % \begin{enumerate}
% %     \item $A\in\cal Q_{2r-1}$ or $B\in \cal Q_{2r-1}$ and $ab\in E(G_t)$,
% %     \item  $A,B\notin \cal Q_{2r-1}$  and $ab\in E(G)$.
% %     %  \item  $B\in \cal Q_{2r}\setminus \cal Q_{2r-1}$ and $A\notin\cal Q_{2r}$  and $ab\in E(G)$.
% % \end{enumerate}


% Let $\cal R$ be the partition of $V(G)-\bigcup \cal Q_{2r}$ which partitions vertices according to their neighborhood in $B_{2r-1}$, in $G$.
% Let $\cal Q$ be the partition of $V$ defined as the union of $\cal Q_{2r}$ and $\cal R$.

% \begin{claim}\label{cl:homogeneous}
%     For every $A\in \cal Q_{2r-1}$ and $B\in \cal R$, the pair $A,B$ is homogeneous in $G$.
% \end{claim}
% \begin{claimproof}
%     Note that for every two vertices $u,v\in B_{2r-1}\cap A$, and for every vertex $w\in B$,
%        $uw\in E(G_t)$ if and only if $vw\in E(G_t)$.
%        Indeed, if $uw\in E(G_t)$ while $vw\notin E(G_t)$, then there is a path of length at most $2r$ from $v$ to $w$ in $(V,R_{t-1})$, contradicting $w\notin \bigcup \cal Q_{2r}$.

%        Since all vertices of $B$ have the same neighborhood in $B_{2r-1}$, in $G$,
%        it follows that all vertices of $B$ are either complete, or anti-complete towards $A$ in $G$.
% \end{claimproof}


% \begin{claim}\label{cl:small-q}
%     $|\cal Q|\le \pi_G(s).$
%    \end{claim}
%    \begin{claimproof}               
%        For every part $A\in \cal Q_{2r-1}$ pick a vertex $v_A\in B_{2r-1}\cap A$,
%        and let $S=\setof{v_A}{A\in\cal Q_{2r-1}}$. In particular, $|S|\le s$.

%        It follows from \Cref{cl:homogeneous} that for every vertex $w\in V(G)-\bigcup \cal Q_{2r}$,
%        the neighborhood  of $w$ in $B_{2r-1}$ is uniquely determined by the neighborhood of $w$ in $S$; moreover, $V(G)-\bigcup \cal Q_{2r}$ is disjoint from $S$.
%        Therefore, $$|\cal R|\le  \pi_G(s)-s.$$
%        Together with $|\cal Q_{2r}|\le s$, this yields the conclusion.
%    \end{claimproof}

% \begin{claim}\label{cl:k-flip'}
%     $G'$ is a $\cal Q$-flip of $G$. In particular, $G'$ is a $\pi_G(s)$-flip of $G$.
% \end{claim}
% \begin{claimproof}
%     It is enough to show that for all   $A,B\in \cal Q$, 
%     \begin{align}
%         E(G')\cap AB=E(G)\cap AB\qquad\textit{or}\qquad
%         E(G')\cap AB= AB-E(G).\label{eq:flip}    
%     \end{align}
%     This is clear if $A,B\in\cal Q_{2r}$, as $G_t$ is a $\cal P_t$-flip of $G$.
%     If one of $A,B$ is in $\cal Q_{2r-1}$ and the other in $\cal R$, then 
%     $A$ and $B$ are homogeneous by \Cref{cl:homogeneous}.
%     As $B\cap B_{2r}=\emptyset$, we have $E(G_t)\cap AB=\emptyset$.
%     It follows that $E(G')\cap AB=E(G_t)\cap AB=\emptyset$.
%     In the remaining cases, by definition of $E(G')$, we have $E(G')\cap AB=E(G)\cap AB$.
% \end{claimproof}



% \begin{claim}\label{cl:balls'}
%     For all $w\in B^r_{G_{t-1}}(v)$
%  and $i=0,\ldots,r$ we have:
%     $$B^i_{G_t'}(w)\subset B^i_{G_t}(w).$$
% \end{claim}

% \begin{claimproof}
%     This is the same as \Cref{cl:balls}.
% \end{claimproof}
% \Cref{cl:k-flip',cl:balls'} together prove \Cref{lem:strategy-poly}.
%     \end{proof}
% \end{szin}

% \begin{szin}\textbf{This proof is fixed (thanks to ChatGPT)}
\begin{proof}
Fix \(v\in V(G)\).  For \(i=0,\ldots,2r\), let
\(
   \mathcal Q_i
\)
be the set of parts \(A\in \cal P_t\) that contain a vertex at distance at most \(i\)
from \(v\) in the graph \((V(G),R_{t-1})\).  
By the assumption on the radius-\(2r\) width of \((\cal P_t,R_{t-1})\), we have
\[
   |\mathcal Q_{2r}|\le s.
\]

Define a graph \(G'=(V(G),E')\) as follows.  For every two parts
\(A,B\in \cal P_t\), set
\[
E'\cap AB =
\begin{cases}
E(G)\cap AB, & \text{if } A,B\notin \mathcal Q_{2r-1},\\
E(G_t)\cap AB, & \text{otherwise.}
\end{cases}
\tag{\(*\)}\label{eq:flip-def}
\]
Note that $G'$ is exactly the same graph as constructed in the proof of \Cref{lem:strategy}. In particular, \Cref{cl:balls} implies that 
$$B^r_{G'}(w)\subset B^r_{G_t}(w)\quad\text{for every $w\in B^r_{G_{t-1}}(v)$.}$$
It therefore remains to show that \(G'\) is a \(\pi_G(s)\)-flip of \(G\); the conclusion of the lemma then follows for $G_t'\coloneqq G'$.

\medskip

For each \(A\in\mathcal Q_{2r-1}\), choose a vertex
\[
   a_A\in A
\]
which is at distance at most \(2r-1\) from \(v\) in \((V(G),R_{t-1})\).
Let
\[
   S:=\{a_A : A\in\mathcal Q_{2r-1}\}.
\]
Then \(|S|=|\mathcal Q_{2r-1}|\le s\).

Denote
\[
   U_{2r}:=\bigcup \mathcal Q_{2r}.
\]
Let \(\mathcal R\) be the partition of \(V(G)\setminus U_{2r}\) in which two
vertices \(y,z\) are in the same part iff
\[
   N_G(y)\cap S = N_G(z)\cap S .
\]
Finally define
\[
   \mathcal Q := \mathcal Q_{2r}\cup \mathcal R .
\]
This is a partition of \(V(G)\).  Moreover,
\[
   |\mathcal Q|
   =
   |\mathcal Q_{2r}|+|\mathcal R|
   \le
   |S|+\bigl|\{N_G(y)\cap S : y\in V(G)\setminus S\}\bigr|
   \le
   \pi_G(s).
\tag{\(**\)}\label{eq:Qdef}
\]

We claim that \(G'\) is a \(\mathcal Q\)-flip of \(G\).  It is enough to show
that for every \(A,B\in\mathcal Q\),
\[
   E'\cap AB = E(G)\cap AB
   \quad\text{or}\quad
   E'\cap AB = AB\setminus E(G).\label{eq:flip-check}
\tag{\(\dagger\)}
\]

If neither \(A\) nor \(B\) contains a part from \(\mathcal Q_{2r-1}\), then
\eqref{eq:flip-def} gives \(E'\cap AB=E(G)\cap AB\).  This handles, in particular, all pairs
inside \(\mathcal R\), and all pairs between \(\mathcal R\) and
\(\mathcal Q_{2r}\setminus\mathcal Q_{2r-1}\).

If \(A,B\in\mathcal Q_{2r}\), then \(A\) and \(B\) are parts of \(\cal P_t\).  On the
pair \(AB\), the graph \(G'\) is either equal to \(G\) or to \(G_t\).  Since
\(G_t\) is a \(\cal P_t\)-flip of \(G\), condition \eqref{eq:flip-check} follows.

It remains to consider the case where \(A\in\mathcal Q_{2r-1}\) and
\(B\in\mathcal R\).  Let \(D\in \cal P_t\) be a part with \(B\cap D\neq\emptyset\).
Since \(B\subseteq V(G)\setminus U_{2r}\), no vertex of \(B\cap D\) is adjacent
to \(a_A\) in \(G_t\).  Indeed, if \(a_Ab\in E(G_t)\) for some \(b\in B\cap D\),
then \(a_Ab\in R_{t-1}\), because \(E(G_t)\subseteq R_t\subseteq R_{t-1}\).
As \(a_A\) is within \(R_{t-1}\)-distance at most \(2r-1\) from \(v\), this would
put the part \(D\) in \(\mathcal Q_{2r}\), contradicting \(B\subseteq
V(G)\setminus U_{2r}\).

Because \(G_t\) is a \(\cal P_t\)-flip of \(G\) and $A,D\in\cal P_t$, it follows that on \(AD\), either
\(G_t\) agrees with \(G\), or with the complement $\overline G$ of \(G\).  The fact that
\(a_A\) has no \(G_t\)-neighbors in \(B\cap D\) implies:
\[
\begin{array}{ll}
G_t \text{ agrees with } G \text{ on } AD
   &\Longrightarrow a_A \text{ is anti-complete to } B\cap D \text{ in } G,\\[2mm]
G_t \text{ agrees with } \overline G \text{ on } AD
   &\Longrightarrow a_A \text{ is complete to } B\cap D \text{ in } G.
\end{array}
\]
But all vertices of \(B\) have the same \(G\)-adjacency to \(a_A\), because
\(B\) is a class of the partition \(\mathcal R\) and \(a_A\in S\).  Hence the
same choice, agreement with $G$ or with $\overline G$, occurs for every \(\cal P_t\)-part
\(D\) meeting \(B\).  Consequently
\[
   E'\cap AB = E(G_t)\cap AB
\]
is either \(E(G)\cap AB\) or \(AB\setminus E(G)\).  This proves
\eqref{eq:flip-check}, and so \(G'\) is a \(\mathcal Q\)-flip of \(G\).  By \eqref{eq:Qdef},
\(G'\) is a \(\pi_G(s)\)-flip of \(G\).
% It remains to prove the ball inclusion.  We show by induction on
% \(i=0,\ldots,r\) that for every \(w\in B^r_{G_{t-1}}(v)\),
% \[
%    B^i_{G'}(w)\subseteq B^i_{G_t}(w).
% \]
% The case \(i=0\) is trivial.  Suppose \(i\ge 1\), and take
% \(u\in B^i_{G'}(w)\).  If \(u=w\), there is nothing to prove.  Otherwise choose
% \(a\in B^{i-1}_{G'}(w)\) with \(au\in E(G')\).  By induction,
% \[
%    a\in B^{i-1}_{G_t}(w).
% \]
% Since \(E(G_t)\subseteq R_t\subseteq R_{t-1}\) and
% \(E(G_{t-1})\subseteq R_{t-1}\), while
% \(w\in B^r_{G_{t-1}}(v)\), we get
% \[
%    a\in B^{r+i-1}_{(V(G),R_{t-1})}(v)
%       \subseteq B^{2r-1}_{(V(G),R_{t-1})}(v).
% \]
% Thus the \(\cal P_t\)-part containing \(a\) belongs to \(\mathcal Q_{2r-1}\).  By the
% definition of \(G'\), every \(G'\)-edge incident with this \(\cal P_t\)-part is an
% edge of \(G_t\).  Hence \(au\in E(G_t)\), and therefore
% \[
%    u\in B^i_{G_t}(w).
% \]
% Taking \(i=r\) gives
% \[
%    B^r_{G'}(w)\subseteq B^r_{G_t}(w)
% \]
% for every \(w\in B^r_{G_{t-1}}(v)\), as required.
\end{proof}

\end{comment}
% \begin{szin}




\begin{lemma}\label{lem:strategy-poly}
Fix \(r\ge 1\) and \(t\in[2,n]\), and let \(s\) be the radius-\((2r-1)\)
width of \((\cal P_t,R_{t-1})\). For every \(v\in V(G)\), there is an
\((s^2+\pi_G(s))\)-flip \(G_t'\) of \(G\) such that
\[
   B^r_{G_t'}(w)\subseteq B^r_{G_t}(w)
   \qquad\text{for every }w\in B^r_{G_{t-1}}(v).
\]
\end{lemma}

\begin{proof}
    Denote $V\coloneqq V(G)$ and 
fix \(v\in V\). Let $\cal N$ 
be the set of parts of \(\cal P_t\) which contain a vertex at distance at most
\(2r-1\) from \(v\) in the graph \((V,R_{t-1})\). 
By the assumption on the radius-\((2r-1)\) width of
\((\cal P_t,R_{t-1})\), we have
\[
   |\cal N|\le s.
\]

Define a graph \(G'=(V,E')\) as follows. For every two parts
\(A,B\in\cal P_t\), set
\[
E'\cap AB =
\begin{cases}
E(G)\cap AB, & \text{if } A,B\notin \cal N,\\
E(G_t)\cap AB, & \text{otherwise.}
\end{cases}
\tag{1}\label{eq:poly-Gprime-def}
\]
Note that $G'$ is exactly the same graph as constructed in the proof of \Cref{lem:strategy}. In particular, \Cref{cl:balls} implies that 
$$B^r_{G'}(w)\subset B^r_{G_t}(w)\quad\text{for every $w\in B^r_{G_{t-1}}(v)$.}$$
It therefore remains to show that \(G'\) is an \((s^2+\pi_G(s))\)-flip of \(G\).



For every \(A\in\cal N\), choose a vertex
\[
   a_A\in A\cap B^{2r-1}_{(V,R_{t-1})}(v),
\]
and put
\[
   S\coloneqq \{a_A:A\in\cal N\}.
\]
Then \(|S|=|\cal N|\le s\).

For \(A\in\cal N\), define
\[
   \mathcal X_A
   \coloneqq
   \{A\}\cup
   \{D\in\cal P_t:N_{G_t}(a_A)\cap D\neq\emptyset\}.
\]
Since \(E(G_t)\subseteq R_t\subseteq R_{t-1}\), every part in
\(\mathcal X_A\) intersects the radius-\(1\) ball of \(a_A\) in
\((V,R_{t-1})\). Since \(1\le 2r-1\), the radius-\((2r-1)\) width gives
\[
   |\mathcal X_A|\le s.
\]
Let
\[
   \mathcal X\coloneqq \bigcup_{A\in\cal N}\mathcal X_A.
\]
Then
\[
   \cal N\subseteq \mathcal X
   \qquad\text{and}\qquad
   |\mathcal X|\le |\cal N|\cdot s\le s^2.
\]

Let
\[
   U\coloneqq \bigcup \mathcal X.
\]
Let \(\mathcal R\) be the partition of \(V\setminus U\) in which two
vertices \(y,z\in V\setminus U\) are equivalent iff
\[
   N_G(y)\cap S=N_G(z)\cap S.
\]
Finally define
\[
   \mathcal Q\coloneqq \mathcal X\cup\mathcal R.
\]
Then \(\mathcal Q\) is a
partition of \(V\). Moreover,
\[
   |\mathcal Q|
   =
   |\mathcal X|+|\mathcal R|
   \le
   s^2+
   \bigl|\{N_G(y)\cap S:y\in V\}\bigr|
   \le
   s^2+\pi_G(s).
\tag{2}\label{eq:poly-Q-size}
\]

We claim that \(G'\) is a \(\mathcal Q\)-flip of \(G\). It is enough to show
that for every \(A,B\in\mathcal Q\),
\[
   E'\cap AB=E(G)\cap AB
   \qquad\text{or}\qquad
   E'\cap AB=AB\setminus E(G).
\tag{3}\label{eq:poly-flip-check}
\]

First suppose \(A,B\in\mathcal X\). Then \(A\) and \(B\) are parts of
\(\cal P_t\). By definition of \(G'\), the graph \(G'\) agrees on \(AB\) either
with \(G\) or with \(G_t\). Since \(G_t\) is a \(\cal P_t\)-flip of \(G\),
\eqref{eq:poly-flip-check} follows.

Next suppose \(A,B\in\mathcal R\). Every \(\cal P_t\)-part meeting \(A\cup B\)
lies outside \(\mathcal X\), hence outside \(\cal N\). Therefore
\eqref{eq:poly-Gprime-def} gives
\[
   E'\cap AB=E(G)\cap AB.
\]

Now suppose \(A\in\mathcal X\setminus\cal N\) and \(B\in\mathcal R\).
Again, every \(\cal P_t\)-part meeting \(B\) lies outside \(\mathcal X\), hence
outside \(\cal N\), and \(A\notin\cal N\). Therefore
\[
   E'\cap AB=E(G)\cap AB.
\]

It remains to consider the case \(A\in\cal N\) and \(B\in\mathcal R\).
Let \(D\in\cal P_t\) be a part with \(D\cap B\neq\emptyset\). Since
\(B\cap U=\emptyset\), we have \(D\notin\mathcal X\). In particular,
\(D\notin\mathcal X_A\), and therefore
\[
   N_{G_t}(a_A)\cap D=\emptyset.
\tag{4}\label{eq:no-Gt-neighbor}
\]
Since \(G_t\) is a \(\cal P_t\)-flip of \(G\), on the pair \(AD\) either
\(G_t\) agrees with \(G\), or \(G_t\) agrees with the complement $\overline G$ of \(G\). By
\eqref{eq:no-Gt-neighbor}, these two alternatives are distinguished by the
\(G\)-adjacency of vertices of \(D\) to \(a_A\):
\[
\begin{array}{ll}
G_t \text{ agrees with } G \text{ on } AD
   &\Longleftrightarrow
   a_A \text{ is anti-complete to } D \text{ in } G,\\[1mm]
G_t \text{ agrees with } \overline G \text{ on } AD
   &\Longleftrightarrow
   a_A \text{ is complete to } D \text{ in } G.
\end{array}
\]
But \(B\) is a part of \(\mathcal R\), so all vertices of \(B\) have the same
\(G\)-adjacency to \(a_A\), because \(a_A\in S\). Hence the same alternative,
agreement with \(G\) or agreement with \(\overline G\), occurs for every
\(\cal P_t\)-part \(D\) meeting \(B\). Since \(A\in\cal N\), we have
\[
   E'\cap AD=E(G_t)\cap AD
\]
for every such \(D\). Therefore \(E'\cap AB\) is either \(E(G)\cap AB\) or
\(AB\setminus E(G)\). This proves \eqref{eq:poly-flip-check}.

Thus \(G'\) is a \(\mathcal Q\)-flip of \(G\). By
\eqref{eq:poly-Q-size}, \(G'\) is an \((s^2+\pi_G(s))\)-flip of \(G\). Taking
\(G_t'\coloneqq G'\) completes the proof.
\end{proof}
% \end{szin}


   As mentioned, \Cref{lem:strategy-poly} implies \Cref{lem:fw-poly},
   analogously as in the proof of \Cref{lem:fw}.


\thmabmw*
\begin{proof}
    Let $\CC$ be a hereditary class of almost bounded merge-width.
    By \Cref{lem:abmw-vc}, $\CC$  has VC-dimension bounded by some $d\in\N$.
By \Cref{lem:sauer-shelah-perles}, we have that $\pi_G(s)\le O(s^d)$ for all $G\in\CC$ and $s\in\N$.

Fix $r\in\N$ and $\eps>0$. Let $\delta>0$ be sufficiently small.
Then, for every graph $G\in\CC$ with $n$ vertices and  $s=\mw_{2r-1}(G)$ 
we have $s\le O_{\CC,r,\delta}(n^\delta)$ and $\pi_G(s)\le O(s^d)$, so by \Cref{lem:fw-poly}:
$$\fw_r(G)\le s^2+\pi_G(s)\le  O(s^{d+2})\le  O_{\CC,r,\delta}(n^{(d+2)\delta}).$$
Setting $\delta=\eps/(d+2)$ gives $\fw_r(G)\le O_{\CC,r,\eps}(n^\eps)$. Thus $\CC$ has almost bounded flip-width.
\end{proof}

% \subsection{Strong Erd{\H o}s-Hajnal}
% Recall that a graph class $\CC$ has the \emph{Strong Erd{\H o}s-Hajnal} property if there is some $\eps>0$ such that every graph $G\in\CC$ contains a pair of vertex subsets $X,Y\subset V(G)$ with $|X|,|Y|\ge >\eps|V(G)|$, and either all or no edges between $X$ and $Y$ in $G$.
% The next lemma shows that graph classes $\CC$ with $\mw_{2}(\CC)<\infty$ have the Strong Erd{\H o}s-Hajnal property. This fact extends the analogous result for twin-width \cite[Thm. 22]{tww3},
% with a very similar proof.

% \begin{lemma}
% The class of graphs of radius-$2$ merge-width at most $w$ satisfies 
% the Strong Erd{\H o}s-Hajnal property with $\eps=1/(w+4)$.
% \end{lemma}
% \begin{proof}Let $G$ be an $n$-vertex graph with $\mw_2(G)=w$.
%   Let $(\cal P_1,R_1),\ldots,(\cal P_n,R_n)$ be a normalized merge-sequence for $G$ of radius-$2$ width $w$ (see \Cref{lem:normalize}).
%   Let $t\in[n]$ be the least such that there is some part $P\in\cal P_t$ with $|P|> n/(d+4)$.
%   Then $|P|\le 2n/(d+4)$, as $P$ is obtained by merging two parts $A,B\in\cal P_{t-1}$ with $|A|,|B|\le n/(d+4)$.

%   Moreover, 

% \end{proof}


% \subsection{Monotone set-flip-width strategies, generalized coloring numbers, and treewidth}


% A \emph{monotone set-flip-width strategy}
% \sz{rewrite}
% is defined like a monotone flip-width strategy, but instead of a partition $\cal P$ of $V(G)$,
% a family $\cal F$ of subsets of $V(G)$ is maintained, which is initially empty.
% Refining the partition is replaced by
% the operation of adding some set $A\subset V(G)$ to the set family $\cal F$,
% while in the second operation, the flip is required to be with respect to some pair of sets in $\cal F$.
% The width is measured in the same way.
% Denote by $\msfw_r(G)$ the minimum radius-$r$ width of any monotone set-flip-width strategy for $G$.


% \begin{lemma}For every graph $G$ and $r\in\N\cup\set{\infty}$, we have:
%     $$\msfw_r(G)\le \mfw_r(G)\le 2^{\msfw_r(G)}.$$
% \end{lemma}
% \begin{proof}
%     The second inequality is immediate.
%     For the first inequality, observe that 
%     a flip between two sets $A,B\in\cal F$
%     can be replaced by a sequence of flips 
%     between pairs of sets belonging to the partition $\cal P$ generated by $\cal F$. 
%     Therefore, the set family $\cal F$ can be replaced by the partition $\cal P$.
% \end{proof}




% For a graph $G$ and a total order $\le$ on $V(G)$, define the \emph{$r$-reachability number}
% as the maximum, over all $u\le v\in V(G)$,
% of the number of vertices $w\le u$ reachable from $v$ by a path of length ${\le}r$,
% whose inner vertices are ${>}u$.

% % \newcommand{\mreach}{\mathrm{reach}}
% \begin{lemma}For every graph $G$ and $r\in\N\cup\set{\infty}$ we have: 
%     $$\scol_r(G) \le \mreach_r(G) \le \wcol_r(G).$$
%     Moreover, $\mreach_\infty(G)=\tw(G)+1$.
% \end{lemma}

% \begin{lemma}
%     \begin{enumerate}
%         \item\label{it:scol} $\msfw_r(G)\le \mreach_{r+1}(G)$,
%          \item\label{it:tw} $\msfw_\infty(G)\le \tw(G)+1$,
%         % \item\label{it:mafw<mw} $\mfw_r\le \mw_r$,
%         % \item\label{it:mafw<mfw} $\mfw_r\le 2^{\msfw_r}$,
%         % \item\label{it:mw<mafw} $\mw_r\le \mfw_r$,
%         % \item\label{it:fw<mafw} $\fw_r\le \mfw_r$.
%     \end{enumerate}
% \end{lemma}
% \begin{proof}
%     \eqref{it:scol}
% Fix a total order $\le$ on $V(G)$ and let $v_1,\ldots,v_n$ be the enumeration of $V(G)$ with respect to that order.
% Define a monotone strategy which performs the following sequence operations, for each $i=1,2,\ldots,n$, maintaining that $E=F=\setof{uv\in E(G)}{u,v\ge v_i}$ 
% and $\cal F=\setof{\set{v_j},\setof{v}{vv_j\in E(G),v>v_j}}{j<i}$
% hold before each step:
% \begin{itemize}
%     \item flip $E$ with respect to $A:=\set{v_i}$ and $B:=N_E(v_i)=\setof{w\in V(G)}{w> v_i, v_iw\in E(G)}$ -- effectively removing all pairs $v_iw\in E$ with $w>v_i$ from $E$ --
%     and add $A,B$ to $\cal F$, 
%     \item next, remove all pairs $v_iw\in F$ with $w>v$ from $F$, ensuring that again $F=E$.
% \end{itemize}
% Clearly, this defines a monotone flip-width strategy for $G$.
% We now bound the radius-$r$ width of this strategy.

% Fix any moment in time $t$ and vertex $v\in V(G)$. We analyse the situation before sequence of operations at time $t$ are performed.
% The ball $B_r^F(v)$ consists of all vertices that are reachable from $v$ by a path of length at most $r$ that does not include any vertex $w$ with  $w<v_i$,
% whereas $\cal F=\setof{\set{v_j},\setof{v}{vv_j\in E(G),v>v_j}}{j<i}$. Hence, the number of sets in $\cal F$ that intersect $B_r^F(v)$ is equal to the number 
% of vertices $w<v_i$ that can be reached from $v$ by a path of length at most $r+1$,
% whose all inner vertices are ${\ge}v_i$.
% This number is exactly $\mreach_{r+1}(G,\le)$.

% \medskip
% \eqref{it:tw}
%     It is known that $\scol_\infty(G)=\tw(G)+1$.
%     Let $r=|V(G)|$. Then $\scol_{r}(G)=\scol_{r+1}(G)=\scol_\infty(G)$ and $\msfw_r(G)=\msfw_\infty(G)$.
%     By the previous item, 
%     $\msfw_r(G)\le \scol_{r+1}(G)$.
%     This proves $\msfw_\infty(G)\le \tw(G)+1$.
% \end{proof}




% \subsection{Treewidth}
% Let $T$ be a tree decompostion of $G$:
% a mapping which associates to each vertex $v$ of $G$  a connected set $T_v$ of nodes of $T$,
% so that for every edge $vw\in E(G)$
% we have $T_v\cap T_w\neq\emptyset$.
% For a node $w$ of $T$,
% the \emph{bag} of $w$ is $\setof{v\in V(G)}{w\in T_v}$, and the \emph{width} of the decomposition is the maximum size of a bag minus one.

% Given a prefix $A$ of $T$,
% define the partition $\cal P_A$ of $V(G)$ so that:
% \begin{itemize}
%     \item All vertices $v\in V(G)$ with nonempty $T_v\cap A$ constitute singleton parts,
%     \item For every maximal node $a\in A$,
%     partition all vertices $v\in V(G)$
%     with $\min T_v>a$ into neighborhood 
%     classes over the bag of $a$. Each of those neighborhood classes is a part of $\cal P_A$.
% \end{itemize}
% Set $R_A\subset E(G)$ as the set of those edges $uv$ of $G$ with $T_u\cap A=T_v\cap A=\emptyset$.

% It follows that every connected component of $R_A$ 

% \begin{claim}
%     Let $uv\in R_A$.
%     Then 
% \end{claim}
% \begin{proof}
    
% \end{proof}



